\chapter{A character family inside a variety of complexes}
\label{ch:character-bridge}

There is a concrete bridge from a complex multiplication order to a
parameterized variety of complexes.  The bridge uses characters of the
underlying rank-two lattice, and its elementary algebra can be written down
completely.  It supplements the continuous flat-torus spectrum with a finite
cochain model, an explicit pullback map, and a quantitative contracting
homotopy.  It also supplies a useful example of cohomology jumping under
specialization.  No cryptographic complexity improvement or interacting
quantum-field spectral estimate is asserted.

\section{A universal complex over the character torus}

Put $S=\mathbb C[X^{\pm1},Y^{\pm1}]$, the Laurent polynomial ring, and
write $A=X-1$, $B=Y-1$.  Consider the cochain complex
\begin{equation}
 0\longrightarrow S\xrightarrow{D_0}S^2
 \xrightarrow{D_1}S\longrightarrow0,
 \qquad
 D_0=\begin{pmatrix}A\\B\end{pmatrix},\quad
 D_1=\begin{pmatrix}-B&A\end{pmatrix}.
 \label{eq:character-universal}
\end{equation}
Its degrees are $0,1,2$.  The identity $D_1D_0=0$ follows from
commutativity.  This is a Koszul complex for two regular parameters, unlike
the periodic resolution in the worked nodal-order example.

\begin{proposition}[Universal exactness and the exceptional fiber]
The universal complex has $H^0=H^1=0$ and
$H^2=S/(X-1,Y-1)\cong\mathbb C$.  After specialization at
$(u,v)\in(\mathbb C^\times)^2$, its differentials are
\[
 d_0=\begin{pmatrix}u-1\\v-1\end{pmatrix},\qquad
 d_1=\begin{pmatrix}-(v-1)&u-1\end{pmatrix}.
\]
The specialized complex is exact if $(u,v)\ne(1,1)$.  At $(1,1)$
its cohomology dimensions are $(1,2,1)$.
\end{proposition}
\begin{proof}
The ring $S$ is a domain, so $D_0$ is injective.  If
$-Bp+Aq=0$, reduction modulo $A$ gives $B\overline p=0$ in
$S/(A)\cong\mathbb C[Y^{\pm1}]$.  The latter ring is a domain and
$B$ is nonzero in it.  Thus $p=Ar$ for some $r\in S$; substitution
and cancellation of $A$ give $q=Br$.  This proves
$\ker D_1=\operatorname{im}D_0$.  The image of $D_1$ is the ideal
$(A,B)$, proving the cokernel assertion.

In a specialized fiber, put $a=u-1$, $b=v-1$.  If $(a,b)\ne(0,0)$,
the first map has rank one, the second is surjective, and their composite
vanishes.  The kernel of the second therefore has dimension one and equals
the image of the first.  If $a=b=0$, both maps vanish and the cohomology is
the entire graded vector space.
\end{proof}

The assignment $(u,v)\mapsto(d_0,d_1)$ is a regular map from an algebraic
torus to the variety of complexes with dimension vector $(1,2,1)$.  Its
image lies in the equation $d_1d_0=0$.  The generic rank pair is $(1,1)$;
the identity character maps to rank pair $(0,0)$.  These are the explicit
rank and determinantal data that connect the example to the general theory
of varieties of complexes \cite{v2BE1973,v2DS1981,v2DEP1982}.

The distinction between the universal cohomology and the fiber cohomology
is essential.  Evaluation at $(1,1)$ is a nonflat base change.  Indeed,
appending $S\to\mathbb C$ to \eqref{eq:character-universal} and
reversing the degree indexing gives a free resolution of the augmentation
module $\mathbb C$.  Tensoring it with that module makes every differential
zero, so
\[
 \operatorname{Tor}_0^S(\mathbb C,\mathbb C)=\mathbb C,\qquad
 \operatorname{Tor}_1^S(\mathbb C,\mathbb C)=\mathbb C^2,\qquad
 \operatorname{Tor}_2^S(\mathbb C,\mathbb C)=\mathbb C,
\]
and higher Tor groups vanish.  The exceptional fiber records derived
information, rather than contradicting the universal exactness calculation.

\section{Characters, local coefficients and sign conventions}

Let $\Lambda=\mathbb Z\omega_1+\mathbb Z\omega_2\subset\mathbb C$,
and $E=\mathbb C/\Lambda$.  A unitary representation
$\rho:\Lambda\to U(1)$ is determined by
$u=\rho(\omega_1)$ and $v=\rho(\omega_2)$.  Here $\rho$ denotes
parallel-transport holonomy.  A bundle with this convention is
\[
 (\mathbb C\times\mathbb C)/
 ((z,w)\sim(z+\lambda,\rho(\lambda)^{-1}w)).
\]
Consequently its sections have automorphy multiplier $\rho^{-1}$.
The flat-torus chapter uses the automorphy character $\chi$;
the conventions agree by setting $\rho=\chi^{-1}$.

A torus has a cell presentation with one vertex, two oriented edges and
one oriented face.  With a consistent choice of lifted cells, its cochains
with these rank-one coefficients have the displayed differences $u-1$
and $v-1$, and the oriented face supplies the row $(-(v-1),u-1)$.
Changing all character conventions to their inverses gives an equivalent
model and leaves the conclusions below unchanged.

The cohomology dimensions also follow directly in the smooth model.
Write $u=e^{2\pi is}$, $v=e^{2\pi it}$, and use real coordinates
$z=x_1\omega_1+x_2\omega_2$.  Coefficient functions have Fourier
frequencies $k=(m-s,n-t)$, with $(m,n)\in\mathbb Z^2$.  On each
mode the flat exterior derivative is exterior multiplication by
$2\pi i(k_1dx_1+k_2dx_2)$.  When $k\ne0$, contraction by
$k_1\partial_{x_1}+k_2\partial_{x_2}$, divided by
$2\pi i|k|^2$, is a contracting homotopy: the exterior multiplication
and contraction anticommutator equals $|k|^2$ times the identity.
For a fixed nontrivial character every mode is nonzero, with frequencies
bounded away from zero.  The resulting multipliers preserve rapidly
decreasing Fourier coefficients and therefore smooth forms.  All
cohomology vanishes.  For the trivial character, only the constant mode
survives and has dimensions $(1,2,1)$.  Thus the finite and differential
complexes exhibit precisely the same cohomology jump.

\section{A complete Hodge calculation}

Equip the three finite cochain spaces with their standard Hermitian inner
products.  For $a=u-1$, $b=v-1$, put
\[
 \sigma(u,v)=|a|^2+|b|^2.
\]

\begin{proposition}[Finite Laplacians and a quantitative homotopy]
The degree-zero, degree-one and degree-two Hodge Laplacians are
\[
 \Delta_0=\sigma I_1,\qquad
 \Delta_1=\sigma I_2,\qquad
 \Delta_2=\sigma I_1.
\]
For $\sigma>0$, the maps $h^1=d_0^*/\sigma$ and
$h^2=d_1^*/\sigma$ satisfy $dh+hd=I$.  Their operator norms are
$\sigma^{-1/2}$.
\end{proposition}
\begin{proof}
The endpoint products $d_0^*d_0$ and $d_1d_1^*$ equal $\sigma$.
At the middle term,
\[
 d_0d_0^*+d_1^*d_1
 =\begin{pmatrix}|a|^2&a\overline b\\b\overline a&|b|^2\end{pmatrix}
 +\begin{pmatrix}|b|^2&-\overline b a\\-\overline a b&|a|^2\end{pmatrix}
 =\sigma I_2.
\]
Division by $\sigma$ gives the homotopy identities in every degree.
Each nonzero row or column differential has operator norm
$\sqrt\sigma$, proving the norm assertion for its adjoint divided by
$\sigma$.
\end{proof}

For unitary characters,
\[
 \sigma=4\sin^2(\pi s)+4\sin^2(\pi t).
\]
Choose $s,t\in[-1/2,1/2]$ and let $\delta^2=s^2+t^2$.  Concavity
of sine on $[0,\pi/2]$ and the inequality $\sin x\le x$ give
\[
 16\delta^2\le\sigma\le4\pi^2\delta^2.
\]
This supplies explicit estimates for the failure of a uniformly bounded
contracting homotopy as the character approaches the identity.  The Euler
characteristic remains $1-2+1=0$ throughout.  Exactness holds at every
nontrivial character, but its quantitative inverse becomes unbounded.
The number $\sigma$ belongs to this chosen finite cellular metric; it is
not the continuum Laplacian eigenvalue formula or a Yang--Mills mass.

\section{An order action with an explicit cochain map}

Suppose $\alpha\Lambda\subseteq\Lambda$ and write its matrix on column
lattice coordinates as
$M=\left(\begin{smallmatrix}a&b\\c&d\end{smallmatrix}\right)$.
Pullback by $[z]\mapsto[\alpha z]$ sends the character pair to
\begin{equation}
 (u',v')=(u^av^c,u^bv^d).
 \label{eq:character-monomial}
\end{equation}
This assertion concerns the multiplier ring of a complex lattice; when
that ring has rank two it is an imaginary quadratic order, as proved in
the flat-torus chapter.

For every integer $r$, define the Laurent polynomial
$S_r(z)=(z^r-1)/(z-1)$, with its polynomial value at $z=1$.
For $r>0$ it is $1+\cdots+z^{r-1}$, for $r=0$ it is zero, and
for $r<0$ it is $-\sum_{j=r}^{-1}z^j$.  In all cases $S_r(1)=r$.
Set
\[
 F_M(u,v)=
 \begin{pmatrix}
 v^cS_a(u)&S_c(v)\\
 v^dS_b(u)&S_d(v)
 \end{pmatrix}.
\]

\begin{theorem}[A character pullback bridge]
The maps $1$, $F_M(u,v)$ and $\det F_M(u,v)$ in degrees $0,1,2$
form a cochain map from the complex at $(u,v)$ to the complex at
$(u',v')$.  At the trivial character the induced maps are
$1$, $M^{\mathsf T}$ and $\det M$, respectively.
If $\alpha\ne0$, precisely $\det M=|\alpha|^2$ unitary characters
pull back to the trivial character.
\end{theorem}
\begin{proof}
The identity
$u^av^c-1=v^cS_a(u)(u-1)+S_c(v)(v-1)$, and its counterpart for
$(b,d)$, show that $d'_0=F_Md_0$.  For any two-by-two matrix $F$ and
column $x$, the alternating row associated to $Fx$, multiplied by $F$,
equals $\det(F)$ times the alternating row associated to $x$.
This follows by expanding the two coordinates, or from
$F^{\mathsf T}JF=(\det F)J$ with
$J=\left(\begin{smallmatrix}0&1\\-1&0\end{smallmatrix}\right)$.
Hence $d'_1F_M=(\det F_M)d_1$, proving the second cochain identity.
At $(1,1)$ the differentials vanish and
$F_M(1,1)=\left(\begin{smallmatrix}a&c\\b&d\end{smallmatrix}\right)
=M^{\mathsf T}$, proving the stated cohomology maps.

Multiplication by a nonzero complex number has real determinant
$|\alpha|^2$.  Changing from a real orthonormal basis to a lattice basis
does not change its determinant, so
$[\Lambda:\alpha\Lambda]=\det M=|\alpha|^2$.
A character pulls back trivially exactly when it is trivial on
$\alpha\Lambda$, equivalently when it factors through the finite group
$\Lambda/\alpha\Lambda$.  A finite abelian group has as many unitary
characters as elements: decompose it into cyclic factors and choose one
root of unity for each generator.  This proves the count.
\end{proof}

The theorem makes an order action, a character pairing, a variety of
complexes, and a cohomology jump coexist in a single explicit diagram.
Its cochain map is an algebraic representative obtained from the displayed
factorization; no choice of a hidden scalar is required.  For example,
$\alpha=n$ kills exactly the $n^2$ torsion characters.  Several initially
acyclic character complexes consequently map to a fiber with nonzero
cohomology.  The map is not asserted to be a cohomology isomorphism.

These finite torsion characters also carry the evaluation pairings used in
finite Fourier analysis.  The present theorem supplies their algebraic
structure and functorial maps.  An efficient encoding, a hidden-function
oracle, and a resource analysis remain separate requirements for a quantum
algorithm.  Similarly, the explicit estimate for $\sigma$ is a bound in
a specified finite complex.  Connecting it to a physical Hamiltonian would
require a proved operator comparison and estimates uniform in the relevant
limits.  The bridge is substantial precisely because its conclusions and
its remaining obligations can both be stated exactly.
